% GENERATED FILE. Edit canonical lessons, then run npm run generate:books.
% canonical-source-hash: fnv1a64:13946c7173a5ea43
% canonical-lessons: SA-C183-yugam, SA-C183-muhurtah, SA-C183-hora, SA-C183-ciram, SA-C183-cirat

\chapter{Age, Auspicious hour, Hour, For a long time, After a long time}
\label{ch:sa-age-auspicious-hour-hour-for-a-long-time-after-a-long-time}
\hlchaptermodality{183}

\begin{chapteropening}
\textbf{By the end of this chapter:} \emph{I can say an age, an auspicious hour, an hour, for a long time and after a long time.}

Complete the last lesson of chapter 183: I can say an age, an auspicious hour, an hour, for a long time and after a long time.
\end{chapteropening}

\section[\texorpdfstring{\sk{युगम्} (yugam)}{युगम् (yugam)}]{\sk{युगम्} (yugam) — an age, an era}

\label{lesson:SA-C183-yugam}



\begin{quote}
Before the new one: say the Sanskrit for gratitude, then the Sanskrit for memory.
\end{quote}

\subsection*{You'll want to know: \sk{युगम्}}
\textbf{\sk{युगम्}} — \emph{yugam} — \textquotedblleft{}an age, an era\textquotedblright{}.

\begin{grammarlens}[title={What you've built}]
One more word for this chapter.
\end{grammarlens}

\subsection*{Your turn}
\begin{itemize}
\raggedright
  \item \emph{Say it:} \emph{yugam}
  \item \emph{Say it:} \emph{yugam}, once more
  \item \emph{Say it:} say \emph{yugam}
  \item \emph{Recall:} say the Sanskrit for gratitude, then the Sanskrit for memory, then say \emph{yugam} again
\end{itemize}

\begin{tcolorbox}[breakable,colback=teal!4,colframe=teal!35!black,title={Before you move on}]
What does \sk{युगम्} mean? (\textquotedblleft{}An age, an era\textquotedblright{}.) Say it once more.
\end{tcolorbox}
\section[\texorpdfstring{\sk{मुहूर्तः} (muhūrtaḥ)}{मुहूर्तः (muhūrtaḥ)}]{\sk{मुहूर्तः} (muhūrtaḥ) — an auspicious hour}

\label{lesson:SA-C183-muhurtah}



\begin{quote}
Before the new one: say the Sanskrit for memory, then the Sanskrit for an age.
\end{quote}

\subsection*{You'll want to know: \sk{मुहूर्तः}}
\textbf{\sk{मुहूर्तः}} — \emph{muhūrtaḥ} — \textquotedblleft{}an auspicious hour\textquotedblright{}.

\begin{grammarlens}[title={What you've built}]
One more word for this chapter.
\end{grammarlens}

\subsection*{Your turn}
\begin{itemize}
\raggedright
  \item \emph{Say it:} \emph{muhūrtaḥ}
  \item \emph{Say it:} \emph{muhūrtaḥ}, once more
  \item \emph{Say it:} say \emph{muhūrtaḥ}
  \item \emph{Recall:} say the Sanskrit for memory, then the Sanskrit for an age, then say \emph{muhūrtaḥ} again
\end{itemize}

\begin{tcolorbox}[breakable,colback=teal!4,colframe=teal!35!black,title={Before you move on}]
What does \sk{मुहूर्तः} mean? (\textquotedblleft{}An auspicious hour\textquotedblright{}.) Say it once more.
\end{tcolorbox}
\section[\texorpdfstring{\sk{होरा} (horā)}{होरा (horā)}]{\sk{होरा} (horā) — an hour}

\label{lesson:SA-C183-hora}



\begin{quote}
Before the new one: say the Sanskrit for an age, then the Sanskrit for an auspicious hour.
\end{quote}

\subsection*{You'll want to know: \sk{होरा}}
\textbf{\sk{होरा}} — \emph{horā} — \textquotedblleft{}an hour\textquotedblright{}.

\begin{grammarlens}[title={What you've built}]
One more word for this chapter.
\end{grammarlens}

\subsection*{Your turn}
\begin{itemize}
\raggedright
  \item \emph{Say it:} \emph{horā}
  \item \emph{Say it:} \emph{horā}, once more
  \item \emph{Say it:} say \emph{horā}
  \item \emph{Recall:} say the Sanskrit for an age, then the Sanskrit for an auspicious hour, then say \emph{horā} again
\end{itemize}

\begin{tcolorbox}[breakable,colback=teal!4,colframe=teal!35!black,title={Before you move on}]
What does \sk{होरा} mean? (\textquotedblleft{}An hour\textquotedblright{}.) Say it once more.
\end{tcolorbox}
\section[\texorpdfstring{\sk{चिरम्} (ciram)}{चिरम् (ciram)}]{\sk{चिरम्} (ciram) — for a long time}

\label{lesson:SA-C183-ciram}



\begin{quote}
Before the new one: say the Sanskrit for an auspicious hour, then the Sanskrit for an hour.
\end{quote}

\subsection*{You'll want to know: \sk{चिरम्}}
\textbf{\sk{चिरम्}} — \emph{ciram} — \textquotedblleft{}for a long time\textquotedblright{}.

\begin{grammarlens}[title={What you've built}]
One more word for this chapter.
\end{grammarlens}

\subsection*{Your turn}
\begin{itemize}
\raggedright
  \item \emph{Say it:} \emph{ciram}
  \item \emph{Say it:} \emph{ciram}, once more
  \item \emph{Say it:} say \emph{ciram}
  \item \emph{Recall:} say the Sanskrit for an auspicious hour, then the Sanskrit for an hour, then say \emph{ciram} again
\end{itemize}

\begin{tcolorbox}[breakable,colback=teal!4,colframe=teal!35!black,title={Before you move on}]
What does \sk{चिरम्} mean? (\textquotedblleft{}For a long time\textquotedblright{}.) Say it once more.
\end{tcolorbox}
\section[\texorpdfstring{\sk{चिरात्} (cirāt)}{चिरात् (cirāt)}]{\sk{चिरात्} (cirāt) — after a long time}

\label{lesson:SA-C183-cirat}



\begin{quote}
Before the new one: say the Sanskrit for an hour, then the Sanskrit for for a long time.
\end{quote}

\subsection*{You'll want to know: \sk{चिरात्}}
\textbf{\sk{चिरात्}} — \emph{cirāt} — \textquotedblleft{}after a long time\textquotedblright{}.

\begin{grammarlens}[title={What you've built}]
One more word for this chapter.
\end{grammarlens}

\subsection*{Your turn}
\begin{itemize}
\raggedright
  \item \emph{Say it:} \emph{cirāt}
  \item \emph{Say it:} \emph{cirāt}, once more
  \item \emph{Say it:} say \emph{cirāt}
  \item \emph{Recall:} say the Sanskrit for an hour, then the Sanskrit for for a long time, then say \emph{cirāt} again
\end{itemize}

\begin{tcolorbox}[breakable,colback=teal!4,colframe=teal!35!black,title={Before you move on}]
What does \sk{चिरात्} mean? (\textquotedblleft{}After a long time\textquotedblright{}.) Say it once more.
\end{tcolorbox}
